% !TEX encoding = UTF-8 Unicode
% !TEX root =  ../Main.tex

\chapter{Einleitung}\label{cha:Einleitung}

% Gliederung in Motivation und Zielsetzung? Und dann den ganzen Theoriekram hier auch rein? 
Verwaltungsprozesse können für Unternehmen und deren Mitarbeitende einen erheblichen Aufwand bedeuten. Dabei treffen häufig veraltete Abläufe auf neue Herausforderungen wie notwendige Effizienz, Transparenz und Datenschutz. Doch selbst bei Unternehmen die nach außen hin modern wirken, können interne Prozesse ineffizient, fehleranfällig und komplex sein. Ferreira da Costa fast das wie folgt zusammen: \glqq{}Eine Webseite ist mit weniger Aufwand modernisiert als etwa ein Bürger-Service oder ein komplexer interner Prozess. Damit liefert eine neue Webseite den Anschein von Modernität, während sie zugleich wenig echten Wandel in der Kommune erfordert (ohne unterschlagen zu wollen, dass die Wartung einer Webseite viel Arbeit in Anspruch nimmt).\grqq{} \NarrativZitat{ferreiradacostaProzessmodernisierungOeffentlichenVerwaltung2023}{2} Dabei ist Prozessoptimierung mithilfe von Softwarelösungen kein Thema, was nur für Technologieunternehmen interessant ist. Auch in Verwaltungen kann durch den Einsatz einer solchen Lösung ein Effizienzgewinn erzielt werden. Neben einer Reduzierung manueller Arbeitsschritte können auch gesetzliche Anforderungen besser eingehalten werden und die Servicequalität für die Zielgruppe deutlich gesteigert werden.  

Ein anschauliches Beispiel liefert der \gls{asta} der \gls{hwr}. Das Sozialbüro des \gls{asta} unter anderem dafür zuständig, Anträge zur Erstattung des Semestertickets zu bearbeiten. Jedes Semester stellen dafür mehrere Hundert Studierende einen Antrag. Es handelt sich hierbei um eine Verwaltungsakt, der zum aktuellen Zeitpunkt weder technisch noch organisatorisch optimiert ist. Die Studierenden reichen ihre Anträge per E-Mail ein und die Daten werden daraufhin händisch in eine Google-Sheets Tabelle übertragen. Dieser Prozess ist nicht nur fehleranfällig sonder auch zeitaufwändig. Hinzu kommt, dass die Speicherung der Daten in Umgebungen erfolgt, die aus Datenschutzsicht fragwürdig sind. In Summe handelt es sich um einen Prozess, der an mehreren Stellen Optimierungspotenziale bietet. 

Das Ziel dieser Arbeit ist es, herauszufinden, welche Schwachstellen der Prozess hat und wie diese mithilfe einer Softwärelösung behoben werden können. Das soll dazu führen, dass der Prozess in Hinsicht auf die Bearbeitungszeit, Fehleranfälligkeit, Nutzerfreundlichkeit (durch die Zielgruppe) und Datenschutzanforderungen optimiert wird. 
Ein Teil der Lösung soll ein Webformular sein, über das Studierende einen Antrag auf der Website des \gls{asta} einreichen können. Den anderen Teil bildet ein internes Bearbeitungsportal, in dem die Mitarbeitenden des Sozialbüros die Anträge einsehen und per Knopfdruck weiterverarbeiten können. Die Arbeit befasst sowohl mit der Erfassung von grundlegenden Anforderungen und den Entwurf der Software als auch mit der technischen Umsetzung. 

Somit ergibt sich die folgende Agenda für die nächsten Kapitel. Zu Beginn werden benötigte Grundlagen zu Datenschutz sowie technischen Kenntnissen erklärt, die für das bessere Verständnis der Arbeit notwendig sind (siehe Kapitel~\ref{cha:Theoretischer_Rahmen}). Im darauf folgenden Kapitel~\ref{cha:Analyse} wird der IST-Prozess beschreiben, daraufhin spezifische Anforderungen an die Software aufgestellt und die Software inhaltlcih als auch optisch modelliert. 

Somit ergibt sich die folgende Agende für die nächsten Kapitel. Zu Beginn werden benötigte Grundlagen zu Datenschutz sowie technischen Kenntnissen erklärt, die für das bessere Verständnis der Arbeit notwendig sind (siehe Kapitel~\ref{cha:Theoretischer_Rahmen}). Im darauf folgenden Kapitel~\ref{cha:Analyse} wird zunächst der aktuelle Prozess analysiert und daraus konkrete Anforderungen an die neue Software abgeleitet. Anschließend wird das Zielbild entwickelt, einschließlich der Modellierung der Benutzeroberflächen, des Soll-Prozesses sowie des zugrunde liegenden Datenmodells. Im Anschluss wird die genaue technische Umsetzung mithilfe des \gls{cms} WordPress beschrieben und abschließend das System in seiner Einsetzbarkeit beim \gls{asta} der \gls{hwr} diskutiert. Dabei steht die Frage im Vordergrund: \glqq{}Wie lässt sich der manuelle Antragsprozess des \gls{asta} der \gls{hwr} durch eine automatisierte Softwarelösung verbessern und kann er so den Datenschutzanforderungen gerecht werden?\grqq{}
